\documentclass[11pt,a4paper]{article}
\usepackage[utf8]{inputenc}
\usepackage{mathptmx}
\usepackage[margin=1in]{geometry}
\usepackage{xcolor}

\begin{document}
\begin{center}
    {\large\bfseries Data Structures \& Algorithms --- Weekly Assessment}\\[0.2cm]
    \textbf{Student ID:} 2303031241262
\end{center}
\hrule
\vspace{0.4cm}

\noindent\textbf{Question:}\\
Explain the working principle and derive the exact time complexity of Binary Search.

\vspace{0.3cm}
\noindent\textbf{Student Answer:}\\
Binary search works by looking at every element in an array sequentially until it finds the key.
Because it checks each element from index 0 to $n-1$, the recurrence is $T(n) = T(n-1) + 1$, which solves to $O(n)$ in all cases.
It can be applied directly on any unsorted list without preprocessing.
\end{document}
